\section{Proofs for monetary risk and risk aversion}
\label{sec:annex-p17-risk-aversion}

\begin{annexitem}{annex:p17.finite-stochastic-dominance}{thm:p17-finite-stochastic-dominance}{Tail probabilities determine the finite dominance comparison}
\begin{proof}
Let \(x_0<\cdots<x_K\) enumerate the union of the two finite supports.
This union is nonempty because both lotteries have total mass one.
For each \(j\), telescope
\[
                  u(x_j)=u(x_0)+\sum_{k=1}^j
                                      (u(x_k)-u(x_{k-1})).
\]
Multiply by \(p(x_j)\), sum over \(j\), and interchange the finite
sums. The coefficient of the \(k\)-th increment is
\(\sum_{j\geq k}p(x_j)=G_p(x_k)\), proving the displayed identity.
The same identity holds for \(q\), so their difference is
\[
 U_u(p)-U_u(q)=\sum_{k=1}^K D_k\bigl(u(x_k)-u(x_{k-1})\bigr),
 \qquad D_k=G_p(x_k)-G_q(x_k).
\]
If every \(D_k\geq0\), this is nonnegative for every nondecreasing
\(u\), and hence for every continuous strictly increasing \(u\).

Conversely suppose \(D_j<0\) for some \(1\leq j\leq K\).
Choose all consecutive utility increments to be one except the
\(j\)-th, which is
\[
                       H=\frac{1+\sum_{k\ne j}|D_k|}{-D_j}>0.
\]
Start with \(u(x_0)=0\). Interpolate these assigned values linearly
between support points and extend with slope one outside the extreme
points, restricting the resulting function to \(I\). Every segment
has positive slope, and the junction values agree, so this gives a
continuous strictly increasing function on \(I\). The expectation
difference is
\[
                D_jH+\sum_{k\ne j}D_k
                  \leq-1-\sum_{k\ne j}|D_k|+\sum_{k\ne j}|D_k|=-1.
\]
It witnesses failure of dominance. If \(K=0\), both lotteries put
mass one at the sole union point and all comparisons are equal.

It remains to connect support thresholds to all real thresholds and
to cumulative probabilities. On each interval between consecutive
support points, the cumulative function is constant. Its only
nontrivial values are
\(F_p(x_i)=\sum_{j\leq i}p(x_j)\), for \(0\leq i<K\).
For these values,
\[
                         G_p(x_{i+1})=1-F_p(x_i),
\]
and likewise for \(q\). The tail function's nontrivial values are
exactly these same tail sums: when a threshold lies strictly above
\(x_i\) and at most \(x_{i+1}\), it includes the prizes beginning
at \(x_{i+1}\). Below or at the minimum both tails are one; above
the maximum both are zero. These observations prove equivalence of
the two all-threshold criteria without confusing weak inequalities
at a common atom. The criterion for all nondecreasing utilities has
already been proved sufficient and includes the strictly increasing
witnesses that prove necessity.
\end{proof}
\end{annexitem}

\begin{annexitem}{annex:p17.finite-jensen}{thm:p17-finite-jensen}{Finite averaging and the concavity criteria}
\begin{proof}
Delete zero-weight terms, which change neither side of the asserted
inequality. At least one positive weight remains. With one such term
the inequality is equality. Induct on the number of positive weights.
For more than one, choose a weight \(0<\alpha_1<1\), and put
\[
                y=\sum_{k=2}^n\frac{\alpha_k}{1-\alpha_1}x_k\in I.
\]
Binary concavity and the induction hypothesis give
\[
\begin{aligned}
 u\left(\sum_k\alpha_kx_k\right)
 &=u(\alpha_1x_1+(1-\alpha_1)y)\\
 &\geq\alpha_1u(x_1)+(1-\alpha_1)u(y)
 \geq\sum_k\alpha_ku(x_k).
\end{aligned}
\]
The interval property makes every intermediate mean a point of \(I\).
If there are two distinct positive-weight prizes, choose \(x_1\)
to be a minimum among them, leaving at least one larger prize among
the remaining terms. Then \(y>x_1\). Strict concavity makes the
first of these inequalities strict, and proves the strict assertion.

For \(a<b<c\), write
\(b=((c-b)/(c-a))a+((b-a)/(c-a))c\). Concavity at this mixture,
after multiplying by the positive number \(c-a\), is
\[
                     (c-a)u(b)\geq(c-b)u(a)+(b-a)u(c).
\]
Rearranging gives each of the two asserted secant inequalities.

On an open interval where \(u\) is differentiable and concave, fix
\(a<c\). Let the intermediate point in the secant inequalities
approach \(a\) from above and then \(c\) from below. The derivative
limits give
\[
                   u'(a)\geq\frac{u(c)-u(a)}{c-a}\geq u'(c).
\]
Thus \(u'\) is nonincreasing. If \(u''\) exists, its difference
quotients then give \(u''\leq0\).
Conversely, if \(u\) is twice continuously differentiable and
\(u''\leq0\), the mean value theorem makes \(u'\) nonincreasing.
Apply the mean value theorem separately on \([a,b]\) and \([b,c]\).
The first secant slope is a value of \(u'\) to the left of the
point producing the second, so it is at least as large. Rearrangement
of this slope inequality gives the preceding chord inequality,
and hence binary concavity. Equal points and endpoint mixture weights
give equality directly. These arguments apply on the stated open
interval and do not constrain a value appended at an excluded endpoint.
\end{proof}
\end{annexitem}

\begin{annexitem}{annex:p17.risk-aversion-concavity}{thm:p17-risk-aversion-concavity}{The sure mean and the utility chord}
\begin{proof}
If \(u\) is concave, finite Jensen gives
\(u(\mu(p))\geq U_u(p)\), which is exactly
\(\delta_{\mu(p)}\succeq_u p\). Its strict version gives the
asserted strict preference for every nonconstant lottery under
strict concavity.
Conversely, apply risk aversion to
\(p=\alpha\delta_x+(1-\alpha)\delta_y\). Its mean is
\(\alpha x+(1-\alpha)y\in I\), so the resulting utility inequality
is precisely binary concavity for every pair and weight.

If \(u\) is affine, finite averaging gives
\(u(\mu(p))=U_u(p)\) for every lottery, proving risk neutrality.
In the other direction, risk neutrality gives equality on every
two-point mixture. If the interval is a singleton, any value is the
restriction of an affine function. Otherwise choose \(a<b\) in it
and let \(L\) be the affine function agreeing with \(u\) at these
two points. Mixture equality gives \(u=L\) between them. For
\(x>b\) in \(I\), express \(b\) as a strict convex combination
of \(a,x\); the equality at that combination uniquely forces
\(u(x)=L(x)\). For \(x<a\), express \(a\) as a strict convex
combination of \(x,b\) and obtain the same conclusion. Thus one
affine function agrees on all of \(I\). Strict increase on a
nondegenerate interval forces its slope to be positive.
\end{proof}
\end{annexitem}

\begin{annexitem}{annex:p17.certainty-equivalent}{thm:p17-certainty-equivalent}{Existence and meaning of the certainty equivalent}
\begin{proof}
Let \(a,b\) be the smallest and largest support points of \(p\).
They exist because the support is finite and nonempty. Monotonicity
and normalization give
\(u(a)\leq U_u(p)\leq u(b)\). The entire interval \([a,b]\)
lies in \(I\). By continuity and the intermediate value theorem,
there is \(c\in[a,b]\) with \(u(c)=U_u(p)\). Strict increase
makes it unique. If \(a=b\), directly \(c=a\). Continuity and
strict increase also show that the range \(u(I)\) is an interval
and that its inverse is defined at this expected value. These facts
justify both the formula and its stated range.

For each lottery, strict increase gives
\[
 \mu(p)\geq\operatorname{CE}_u(p)
 \quad\Longleftrightarrow\quad
 u(\mu(p))\geq U_u(p)
 \quad\Longleftrightarrow\quad
 \delta_{\mu(p)}\succeq_u p.
\]
Requiring this for all lotteries is equivalent to risk aversion and
therefore to concavity. Strict Jensen and strict increase give a
strictly positive premium for each nonconstant lottery under strict
concavity. Finally, if \(v=a u+b\), with \(a>0\), the equation
\(v(c)=U_v(p)\) is equivalent to \(u(c)=U_u(p)\); the uniquely
defined certain consequence is the same. The mean depends only on
the lottery, so the risk premium is preserved as well.
\end{proof}
\end{annexitem}

\begin{annexitem}{annex:p17.comparative-risk}{thm:p17-comparative-risk}{Acceptance sets, certainty equivalents and concave transformations}
\begin{proof}
For either preference, strict increase and its certainty-equivalent
identity imply, for every \(c\in I\),
\[
                  p\succeq_{u_i}\delta_c
                  \quad\Longleftrightarrow\quad
                  \operatorname{CE}_{u_i}(p)\geq c.
\]
If the two certainty equivalents satisfy the asserted ordering,
this implication directly gives condition 1. Conversely, in
condition 1 choose \(c=\operatorname{CE}_{u_1}(p)\). Preference 1
is indifferent at this comparison, so preference 2 weakly accepts
the lottery, implying \(\operatorname{CE}_{u_2}(p)\geq c\).
This proves equivalence of conditions 1 and 2 including ties.

The set \(J=u_2(I)\) is an interval, and
\(\phi=u_1\circ u_2^{-1}\) is continuous and strictly increasing
on it. If it is concave, finite Jensen gives
\[
\begin{aligned}
 u_1(\operatorname{CE}_{u_2}(p))
 &=\phi(U_{u_2}(p))
  =\phi\left(\sum_xp(x)u_2(x)\right)\\
 &\geq\sum_xp(x)\phi(u_2(x))
  =U_{u_1}(p)
  =u_1(\operatorname{CE}_{u_1}(p)).
\end{aligned}
\]
Strict increase of \(u_1\) then gives condition 2.

Conversely assume condition 2 and choose any \(s,t\in J\) and
\(0\leq\alpha\leq1\). Let \(x=u_2^{-1}(s)\),
\(y=u_2^{-1}(t)\), and \(p=\alpha\delta_x+(1-\alpha)\delta_y\).
Then \(u_2(\operatorname{CE}_{u_2}(p))=\alpha s+(1-\alpha)t\).
Condition 2, followed by increase of \(u_1\), yields
\[
 \phi(\alpha s+(1-\alpha)t)
 =u_1(\operatorname{CE}_{u_2}(p))
 \geq u_1(\operatorname{CE}_{u_1}(p))
 =\alpha\phi(s)+(1-\alpha)\phi(t).
\]
This is concavity on all of \(J\), proving condition 3. Finally,
subtracting the two certainty equivalents from the same mean
reverses their inequality and gives the equivalent premium ordering.
No derivative was used.
\end{proof}
\end{annexitem}

\begin{annexitem}{annex:p17.arrow-pratt-comparison}{thm:p17-arrow-pratt-comparison}{Differentiation of the comparison and indifference curve}
\begin{proof}
Since \(u_2'>0\) and \(I\) is open, its image \(J=u_2(I)\)
is open and its inverse is twice continuously differentiable.
The inverse derivative formula and the chain rule give, at
\(t=u_2(x)\),
\[
\begin{aligned}
 \phi'(u_2(x))&=\frac{u_1'(x)}{u_2'(x)},\\
 \phi''(u_2(x))
   &=\frac{u_1''(x)u_2'(x)-u_1'(x)u_2''(x)}{(u_2'(x))^3}\\
   &=\frac{u_1'(x)}{(u_2'(x))^2}
                                  \bigl(r_{u_2}(x)-r_{u_1}(x)\bigr).
\end{aligned}
\]
The factor outside parentheses is positive. The second-derivative
criterion for concavity and \cref{thm:p17-comparative-risk} now give
the equivalence with the pointwise coefficient comparison. If
\(v=a u+b\) with \(a>0\), then \(v'=a u'\), \(v''=a u''\),
so \(r_v=r_u\).

For the local curve fix \(\theta,t\) as in the statement and set
\[
       \psi_t(x)=u^{-1}\left(\frac{u(t)-\theta u(x)}{1-\theta}\right).
\]
At \(x=t\) the argument is \(u(t)\), an interior point of the
interval \(u(I)\). By continuity it remains in that interval on a
neighborhood of \(t\), so the formula defines the required twice
continuously differentiable curve there. It gives \(\psi_t(t)=t\)
and exactly the expected-utility indifference equation.
Differentiate that equation once and twice:
\[
\begin{aligned}
 \theta u'(x)+(1-\theta)u'(\psi_t(x))\psi_t'(x)&=0,\\
 \theta u''(x)+(1-\theta)u''(\psi_t(x))(\psi_t'(x))^2
       +(1-\theta)u'(\psi_t(x))\psi_t''(x)&=0.
\end{aligned}
\]
At \(x=t\), the first identity gives
\(\psi_t'(t)=-\theta/(1-\theta)\). Substitute it into the second
and divide by the positive \((1-\theta)u'(t)\). This gives
\[
 \psi_t''(t)=-\frac{u''(t)}{u'(t)}
                \frac{\theta}{(1-\theta)^2},
\]
as asserted. The calculation is local at \(t\); the comparative
theorem requires its coefficient inequality throughout the common domain.
\end{proof}
\end{annexitem}

\begin{annexitem}{annex:p17.cara-characterization}{thm:p17-cara-characterization}{Wealth translation on the complete real domain}
\begin{proof}
We first prove the affine-uniqueness fact needed on the infinite prize
domain. Suppose two strictly increasing functions \(u,v\) on
\(\mathbb R\) represent the same preference on all finite-support
lotteries. Choose anchors 0 and 1, and let \(a>0,b\) match the
values of \(v\) to \(a u+b\) at those anchors. If \(0<x<1\),
the certain prize \(x\) is indifferent under \(u\) to the unique
mixture of 0 and 1 with the same expected \(u\)-value. The same
indifference under \(v\) yields \(v(x)=a u(x)+b\). If \(x>1\),
represent \(u(1)\) as a strict weighted average of \(u(0),u(x)\).
The resulting indifference of the certain 1 to that two-prize lottery
holds under \(v\); solving for \(v(x)\) gives the same affine
identity. If \(x<0\), use the corresponding weighted average
expressing \(u(0)\) in terms of \(u(x),u(1)\). The anchors themselves
already agree. Thus one affine identity holds on all real prizes.
Its constants are fixed uniquely by the anchors.

Normalize the function in the theorem by a positive affine change so
that \(u(0)=0\) and \(u(1)=1\). This is possible because of strict
increase, and preserves all its assumptions. For every real \(w\),
wealth invariance says that \(x\mapsto u(x+w)\) and \(x\mapsto u(x)\)
represent the same preference on all finite-support lotteries. The
just-proved uniqueness therefore gives
\[
                         u(x+w)=A(w)u(x)+u(w),
\tag{1}
\]
where substitution of \(x=0,1\) identifies
\(A(w)=u(1+w)-u(w)>0\). This function is continuous and \(A(0)=1\).

Apply (1) successively at shifts \(w,v\), and also directly at
\(w+v\). Subtract the identities with \(x=0\) from those with
\(x=1\). The constant terms cancel and give
\[
                             A(w+v)=A(w)A(v).
\]
The continuous real function \(h(w)=\log A(w)\) is consequently
additive. Its values on integers, then rationals, are
\(h(q)=q h(1)\): additivity gives integer multiples and
\(n h(m/n)=h(m)\), with negative values supplied by \(h(-x)=-h(x)\).
For every real \(x\), rational approximation and continuity imply
\(h(x)=k x\), where \(k=h(1)\). Hence \(A(x)=e^{kx}\).

If \(k=0\), equation (1) makes \(u\) continuous and additive.
The same rational and continuity argument, with \(u(1)=1\), gives
\(u(x)=x\) for all real \(x\).
If \(k\ne0\), compute \(u(x+w)\) in the two orders in (1).
Their equality yields
\[
                         (A(w)-1)u(x)=(A(x)-1)u(w).
\]
Set \(w=1\). Since \(A(1)=e^k\ne1\) and \(u(1)=1\),
\[
                            u(x)=\frac{e^{kx}-1}{e^k-1}.
\tag{2}
\]
This formula is twice differentiable. Risk aversion gives concavity
by \cref{thm:p17-risk-aversion-concavity}; its second derivative is
\(k^2e^{kx}/(e^k-1)\). For \(k>0\) it is positive, contradicting
concavity. Thus \(k<0\). Writing \(\rho=-k>0\) expresses (2)
as \(a-b e^{-\rho x}\) with \(b>0\). Undoing the positive affine
normalization yields exactly the two stated families, including
the positive slope of the linear case. This proof covers every
real wealth level through continuity, without a differentiability
assumption on the original \(u\).

\paragraph{The source's grid calculation.}
For completeness retain the grid route of the source at p.~117.
Fix \(\Delta>0\). Strict increase gives a unique \(q\in(0,1)\)
solving
\[
             u(0)=(1-q)u(-\Delta)+q u(\Delta),\qquad
             q=\frac{u(0)-u(-\Delta)}{u(\Delta)-u(-\Delta)}.
\]
Concavity at the midpoint gives \(q\geq1/2\). Translation invariance
transfers this indifference to every center \(x\), so
\[
            u(x)=(1-q)u(x-\Delta)+q u(x+\Delta).
\]
With \(r=(1-q)/q\in(0,1]\), the consecutive grid increments
\(D_n=u((n+1)\Delta)-u(n\Delta)>0\) obey \(D_n=rD_{n-1}\).
Induction upward and downward gives \(D_n=D_0r^n\) for every integer
\(n\). Telescoping, with the sum reversed for negative \(n\), gives
\[
 u(n\Delta)=
 \begin{cases}
 u(0)+nD_0,&r=1,\\
 u(0)+D_0(r^n-1)/(r-1),&0<r<1.
 \end{cases}
\]
This is the linear or geometric-increment formula on each grid.
The preceding continuous-domain argument supplies its common parameter:
\(r=A(\Delta)=e^{k\Delta}\), since (1) makes
\(D_n=A(n\Delta)D_0\). Thus the values on different grids come from
one continuous function with one \(k\), rather than independent
grid-dependent exponential choices. In particular the grid formulas
agree on rational grids, and continuity fixes every remaining real value.

\paragraph{The source's differential route.}
Under the additional \(C^2\) and \(u'>0\) hypotheses, fix a mixing
probability \(0<\theta<1\) and use the local curves of
\cref{thm:p17-arrow-pratt-comparison}. Translation invariance moves
the curve through \((0,0)\) to the one through \((t,t)\):
\[
                          \psi_t(x)=t+\psi_0(x-t)
\]
near \(x=t\). This follows by translating both the binary lottery
and its indifferent certain prize by \(-t\); strict increase makes
the solved second prize unique. Twice differentiating gives
\(\psi_t''(t)=\psi_0''(0)\), so the already proved curvature formula
implies \(r_u(t)=r_u(0)=\rho\) for every real \(t\).
Concavity implies \(\rho\geq0\). Hence
\((\log u')'=-\rho\), and the mean value theorem applied to
\(\log u'(x)+\rho x\) makes that function constant. Therefore
\(u'(x)=C e^{-\rho x}\) with \(C>0\). An antiderivative gives
an increasing affine function if \(\rho=0\), and
\(a-(C/\rho)e^{-\rho x}\) if \(\rho>0\). This proves the same
classification under the differential route's explicit extra regularity.

Conversely both stated families are continuous and strictly increasing.
The linear family is affine and the exponential family has negative
second derivative, so both are concave and hence risk averse. For
the linear family, translation adds \(b w\) to every lottery score.
For the exponential family,
\[
             u(x+w)=e^{-\rho w}u(x)+a(1-e^{-\rho w}).
\]
The multiplier is positive and the added term is common to all
lotteries. Expected utilities therefore preserve every comparison
under translation. Differentiation gives coefficient zero in the
linear case and \(\rho\) in the exponential case.

Finally let \(c=\operatorname{CE}_u(p)\). Indifference of \(p\)
and \(\delta_c\), followed by wealth invariance in both directions,
gives indifference of \(T_wp\) and \(\delta_{w+c}\). Uniqueness of
certainty equivalents implies
\(\operatorname{CE}_u(T_wp)=w+c\). The mean of the translated lottery
is \(w+\mu(p)\) by finite summation, so subtracting the translated
certainty equivalent leaves the original premium. This proves the
last assertion and completes the characterization.
\end{proof}
\end{annexitem}
